% =====================================================================
% PROPOSTA -- non incluso in main.tex
%
% Scopo: chiudere la seconda meta' del punto 3 della consegna,
%   "The coursework must discuss how structural aliasing at the signal and
%    token levels may affect the correct recognition of these domain-level
%    concepts."
%
% Dove va: in coda a sections/three.tex, dopo il paragrafo che finisce con
%   "Confirmed aliasing suspends forecast-driven battery decisions."
%   Quel paragrafo va sostituito dal blocco qui sotto, che lo riassorbe
%   e ne recupera le due ancore perse nell'accorciamento:
%     - il rimando a \autoref{tab:bayesDecisions} per la soglia di conferma;
%     - la conseguenza per il caso intermedio (ConservativeMode).
%
% Costo: ~250 parole. Margine disponibile: ~470 (corpo a 8327 su 8800).
%
% Vocabolario verificato sull'Appendice A (versione corrente):
%   Anomaly            : InverterTrip, RampEvent, SensorMalfunction
%   Generation context : ClearSky, PartialCloud, Overcast, Nighttime
%   Observability      : Observable, PartiallyObservable, StructurallyAliased
%   Control mode       : NormalMode, ConservativeMode, SafeMode
% =====================================================================

Geometry alone yields a candidate, never a verdict. Because the confirmation
threshold of \autoref{tab:bayesDecisions} exceeds one half, confirmed loss and
confirmed preservation cannot both hold. A frequency for which preservation is
positively confirmed is assigned \valueof{Observable}; everything else is
assigned \valueof{PartiallyObservable}: an unfinished test, a credible interval
crossing the threshold, or conflicting metrics.

\subsection{How aliasing reaches the domain concepts}

The concepts of the ontology are not read off the raw telemetry. A
\valueof{RampEvent}, a \valueof{PartialCloud} generation context or an
\valueof{InverterTrip} is inferred from an observation whose short-horizon
forecast the agent has already trusted, so a frequency the representation
cannot hold is a frequency whose signature cannot support that inference. The
exposure is therefore selective, and the arithmetic says which concepts it
selects. On the one-minute telemetry of this domain, with $P=16$, the patch
comb sits at periods of $16/k$ minutes, that is at $16$, $8$, $5.3$, $4$,
$3.2$, $2.7$ and $2.3$ minutes. Those are the timescales of cloud-edge
crossings and of inverter duty cycling, so the fluctuation that distinguishes
\valueof{PartialCloud} from \valueof{Overcast}, and the leading edge by which a
\valueof{RampEvent} is recognised, fall inside the candidate set. A
\valueof{SensorMalfunction} or an \valueof{InverterTrip} does not: a step is
broadband, and no single comb can hide it.

The failure this predicts is not a wrong number but a missed distinction. Two
generation contexts that differ only in content at a candidate frequency are
mapped towards the same latent state, so the state is silently misassigned
rather than flagged as uncertain, and the anomaly whose onset carries that
frequency is recognised late or not at all. That is precisely why the
assessment, and not the geometry, selects the control mode. Confirmed aliasing
forces \valueof{SafeMode} and suspends forecast-driven battery decisions; a
merely \valueof{PartiallyObservable} assessment forces \valueof{ConservativeMode},
which reduces the battery cap and lowers decision confidence. The inconclusive
verdict of the Bayesian analysis is thus carried into the semantics instead of
being resolved by assumption, and uncertainty about the representation reaches
the actuator rather than stopping at the report.
